\section*{Results}

% Owner: papers/_program/program.md §1 (body roster and figure order) and
% papers/1_method/decisions.md ("The figure set", revised 2026-07-28: six body figures).
% Every quoted number is on the freeze (1c2ae6f + 87889e6 + 0ffbda7), anchored on the Gaussian
% Fisher information, at n_sim = 10^4 unless the sentence says otherwise, and carries a `% src:`
% comment per LINT-SRC. Nothing is quoted from 433ed13.
% Definitions belong to 02_theory.tex and 03_diagnostics.tex; run values, the two least-squares
% configurations and the noise conversion belong to 06_methods.tex. Do not restate them here.

The figures are ordered so that each one answers the objection the one before it raises. Figure~\ref{fig:ladder} shows that the members of the cost ladder do different arithmetic on the same recording, which invites the question of whether that arithmetic changes what an experimenter takes away. Figure~\ref{fig:clouds} answers it at one design cell, where the members recover the parameters and report intervals that differ by more than an order of magnitude. That raises the question of whether the fault lies in the likelihood or in the fit, so Figure~\ref{fig:time} puts the question to the likelihood itself, scoring one common set of recordings interval by interval, and locates the failure in how the information accumulates over time. One cell is not a result, so Figure~\ref{fig:plane} repeats both moments over the plane of channel number against instrumental noise, with the acquisition interval resolved inside each cell. A member whose reported interval is honest over most of that plane invites the question of where it stops being so, and Figure~\ref{fig:design} answers it on that corner, decomposing the departure into the two discarded moments that cause it. Figure~\ref{fig:usage} collects the answers into a map of what a macroscopic recording can be asked for.

Throughout, $N_{\mathrm{ch}}$ is the number of channels in the patch, $k_{\mathrm{on}}$ and $k_{\mathrm{off}}$ the opening and closing rates, $i$ the unitary current, $\Delta$ the acquisition interval and $\tau = 1/k_{\mathrm{off}}$ the closing time constant, so that $\Delta\!\cdot\!k_{\mathrm{off}}$ is the interval in units of $\tau$. The instrumental noise is reported as the label the runs and the figure axes carry; the dimensionless noise $\tilde S$, the instrumental noise power accumulated over one channel time constant measured against the unitary current, is one tenth of that label (Methods).

\subsection*{One filter step along the cost ladder}

Figure~\ref{fig:ladder} puts one turn of the filter cycle for the four body members side by side, on a single simulated recording of a patch of $N_{\mathrm{ch}}=20$ channels. It fixes the vocabulary and carries no measured claim.

Least squares (\texttt{LSE}) fits the deterministic mean current and assigns one constant variance $\hat\sigma^2$ to every residual, estimated as the residual sum of squares over the record. Its predicted band therefore has the same width at every interval, whether the channels are gating or not, which is the homoscedastic premise of the classical fit made visible on the trace. The non-recursive member (\texttt{NR}) replaces that single constant with the per-interval gating variance the model predicts, so its band follows the open population, and its update row reads as an open loop: the occupancy carried into each interval comes from the initial condition and never from the data. The recursive member (\texttt{R}) adds the Bayesian update, on a conductance taken at a single instant. The boundary-conditioned member (\texttt{IR}) conditions the interval-averaged conductance on the channel states at both ends of the interval, which lets the covariance between them enter the gain. Reading left to right, each column adds one operation to the column before it, and the cost of a likelihood evaluation rises with it.

\begin{figure}[htbp]
\begin{fullwidth}
\centering
\includegraphics[width=\linewidth]{Figure_1}
\end{fullwidth}
\caption{\textbf{One filter step along the cost ladder.} The current from a patch of $N_{\mathrm{ch}}=20$ two-state channels, recorded as an interval average with white instrumental noise, is scored by each of the four body members (columns): least squares on the deterministic mean current (\texttt{LSE}), the open-loop gating likelihood (\texttt{NR}), the recursive filter on instantaneous samples (\texttt{R}), and the boundary-conditioned filter (\texttt{IR}). Rows follow one turn of the cycle: (A) the prior open probability from the Markov step, (B) the predicted current with its spread and the innovation, (C) the Bayesian posterior, which for the two non-recursive columns reads as no update, and (D) the per-step log-likelihood accumulated over the window and reset at the window start. Two measurement windows are shown (6 to 10 ms of a 12 ms recording). \texttt{LSE} is drawn in \texttt{NR}'s grammar, a point mean with an error bar, so that its constant least-squares variance can be read against \texttt{NR}'s per-interval gating variance. Colour encodes the filter phase and is shared across rows: green the simulated channel current, orange the predict phase, blue the observation and the innovation, purple the update, black the log-likelihood. A quantity carried in from the neighbouring interval is drawn hollow or dashed. Three of the seven rungs are not drawn, because the cycle shown here does not change for them: \texttt{VR} differs from \texttt{MR} in the interval variance alone and their gains are identical, while \texttt{INR} and \texttt{MR} differ from \texttt{NR} and \texttt{R} only in which conductance vector enters, which moves the predicted moments and the gain built from them without altering the form of any step. The one step that does change form, the gain acquiring the boundary term, is the \texttt{R}-to-\texttt{IR} step drawn here. Appendix~1 writes all seven out. The recording is one $20$-channel trace simulated with $250$ within-interval substeps and does not enter any quantitative comparison.} % src: projects/eLife_2025/figures/paper_both/figure_1.Rmd (merged roster LSE/NR/R/IR; LSE variance = SSE/n) ; papers/_program/provenance.md#4 (20 channels, 250 substeps)
\label{fig:ladder}
% NOTE: the graphics file is still named Figure_S1.pdf on disk, from the retired independent
% supplementary numbering. It is written by figure_1.Rmd on the merged roster and is current;
% rename it to Figure_1_supplement_1.pdf in the deposit pass.
\figsupp{The same filter cycle over the first ten milliseconds of the recording, so that the concentration jump and the onset of gating are inside the window.}{\includegraphics[width=\linewidth]{Figure_S1}}
\end{figure}

\subsection*{Recovery at one design cell}

% Rebuilt 2026-07-30 on the full seven-rung roster and three parameter pairs; the IR correlation
% corner that used to be panel B moved out of the body figure. Numbers below are the figure's own,
% recomputed from figure_2.Rmd's diagnostics chunk.
Figure~\ref{fig:clouds} turns the ladder into a statement about reported uncertainty at one cell: $N_{\mathrm{ch}}=100$, noise label $0.1$ ($\tilde S = 0.01$), $\Delta\!\cdot\!k_{\mathrm{off}} = 0.1$. Ten thousand independent recordings are partitioned into groups of ten, each group is fitted by maximum likelihood under each member, and the scatter of the resulting thousand estimates is compared against the covariance that member reports from its own Gaussian Fisher information.

The two objects disagree by an order of magnitude at the foot of the ladder and by nothing at its top, but the fall between them is not monotone.
% TODO-SIGN: direction convention as fixed in Diagnostics; if the pending code read reverses it,
% flip every TODO-SIGN site in one pass.
Taking the ratio of the area of the empirical $95\%$ ellipse to the area of the reported one, and averaging geometrically over the six pairs the four informative parameters form, the open-loop member reads $11.41$ (per-pair range $10.04$ to $14.83$), the recursive instantaneous member $1.17$ ($1.08$ to $1.32$), and the boundary-conditioned member $1.01$ ($0.98$ to $1.05$). % src: papers/1_method/decisions/D-4_ranking_verdict.md §1a (recompute/d4_headline_and_nch_trend.py, freeze + Gaussian anchor)
On the joint four-parameter ellipse the same three read $154.65$, $1.56$ and $1.12$. % src: papers/1_method/decisions/D-4_ranking_verdict.md §1a
Per parameter, the diagonal of the distortion matrix $\mathbf{C}$ at this cell runs $13.24 / 17.32 / 10.04 / 9.91$ for \texttt{NR}, $1.39 / 1.51 / 1.16 / 1.16$ for \texttt{R} and $1.05 / 1.13 / 1.00 / 0.90$ for \texttt{IR}, on $k_{\mathrm{on}}$, $k_{\mathrm{off}}$, $i$ and $N_{\mathrm{ch}}$ in that order. % src: papers/1_method/decisions/D-4_ranking_verdict.md §1a
A likelihood that carries no recursion therefore under-reports the width of the parameter distribution by a factor of ten to fifteen, while the recursive members are wrong by tens of percent. % src: papers/1_method/decisions/D-4_ranking_verdict.md §3 (10 to 15, ellipse-area factor for NR at this cell)
The one-endpoint members do not sit between the two recursive ones, and this is where the ladder stops being an ordering. On the kinetic pair the recursive instantaneous filter reads $1.32$ while \texttt{MR} reads $1.97$ and \texttt{VR} reads $2.19$, and only then does \texttt{IR} return to $1.02$; the amplitude pair runs $1.09$, $1.53$, $1.77$, $1.00$ and the noise pair $1.06$, $1.29$, $1.43$, $0.99$ in the same order. % src: projects/eLife_2025/figures/paper_both/figure_2.Rmd (diagnostics chunk, empirical over Fisher ellipse area, three pairs)
Conditioning the conductance on one endpoint therefore costs calibration rather than buying it, in every pair and in both directions of the variance switch, and what recovers it is the second endpoint entering the gain (Appendix~1). The two intermediate rungs are kept in the figure for exactly this reason: they are the step from \texttt{R} to \texttt{IR} taken in its two algebraic halves, and the half that changes only the interval variance moves the answer the wrong way.

What the acquisition average buys on its own is separable here, because one member isolates it: the one-endpoint non-recursive member conditions the conductance on the interval instead of on an instant and adds nothing else, no update and no second endpoint. Against the open-loop member it improves every pair, from $14.83$ to $12.59$ on the kinetic pair and from $10.04$ to $9.00$ on the amplitude pair, and on the noise pair it improves by nearly a factor of three, from $8.84$ to $3.28$. % src: projects/eLife_2025/figures/paper_both/figure_2.Rmd (diagnostics chunk, empirical over Fisher; macro_INR from 1f7138b)
The reason is in the first moment rather than the second. The open-loop member carries no interval-mean conductance variance, so the gating variance it predicts is too small, and the fit makes up the difference by inflating the one free parameter that can absorb it: its noise estimate lands a factor of $2.2$ above the truth. Restoring that term removes the need to compensate, and the noise estimate returns to within $0.005$ in $\log_{10}$. Modelling the acquisition average therefore buys the noise level, and it buys it with no recursion at all.

The distortion is measurable and it is also removable. Replacing each reported covariance by the sandwich, which substitutes the observed covariance of the score for the information the likelihood claimed, returns every gating-aware member of the ladder to within about a tenth of one: $1.03$ for \texttt{NR}, $0.95$ for \texttt{R}, $1.05$ for \texttt{IR}, with every parameter pair falling between $0.87$ and $1.07$, and the two one-endpoint members inside the same band. % src: papers/1_method/decisions/D-4_ranking_verdict.md §1a (emp/corrected column) ; figure_2.Rmd diagnostics chunk (MR 0.88-0.93, VR 0.94-0.98 over the three pairs)
The correction repairs the members that model the fluctuations, whatever they get wrong about them, and it is least squares alone that it leaves short.
The reader-facing form of that repair is coverage rather than area. Reducing the full multivariate distribution of the estimates to the squared Mahalanobis distance from the cloud mean, which is distributed as $\chi^2$ with as many degrees of freedom as free parameters if the reported covariance is right, the nominal $95\%$ region of \texttt{IR} covers $0.914$ of the estimates under its own Fisher covariance and $0.947$ to $0.949$ under the sandwich, measured at $N_{\mathrm{ch}}=10$ and noise label $0.05$ over a thousand fits in six parameters at group size $100$ (Figure~\ref{fig:plane}--figure supplement 4). % src: projects/eLife_2025/figures/paper_both/figure_4_supplement_5.html (IR, N_ch 10, noise 0.05, 1000 fits at group size 100, p = 6, intervals 1 and 0.05)
That one correction repairs six of the seven rungs, without being told which of them is wrong or how, is what makes the distortion a usable quantity rather than a complaint.

Least squares carries a column in the kinetic pair and none in the other two, and the reason is identifiability rather than a gap in the runs. Its grid configuration holds the unitary current and the instrumental noise level fixed at the values the recordings were simulated from (Methods), so in the amplitude pair its estimates would lie exactly on the line $i=1$ and in the noise pair exactly on $\sigma = 10^{-4}$: the empirical and the reported ellipse would both have zero area and the ratio the panel reports would be $0/0$. The mean current depends on the channel number and the unitary current only through their product, and it is the gating fluctuations, whose variance carries a second power of the unitary current, that separate the two. A method that models no fluctuations estimates the product, and fixing the unitary current names that product after one of its factors. Where least squares does compete, on the pair the deterministic transient determines, it under-reports by $14.8$, which is the open-loop member's own figure, and it is the only member of the seven the sandwich does not repair: corrected, it still reads $1.63$ against $0.96$ to $1.07$ for every other member and pair. % src: projects/eLife_2025/figures/paper_both/figure_2.Rmd (diagnostics chunk: LSE on/off 14.78 empirical over Fisher, 1.63 over corrected)

\begin{figure}[htbp]
\begin{fullwidth}
\centering
% Rebuilt 2026-07-30 from figures/paper_both/figure_2.Rmd on the full seven-rung roster. The three
% members that were missing resolve from three different engine-commit directories (LSE and VR from
% RESOLVED 2026-08-01: the third column is macro_INR, from the corrected engine (1f7138b, clean
% stamp), run on dirac for this cell alone. It replaced macro_NMR, the build that had lost the N*ms
% term since a3e0a89; the two are different algorithms and not two spellings (nomenclature.md).
% Swapping them changed the column: NMR read 14.0 / 9.8 / 9.0 on the three pairs and INR reads
% 12.6 / 9.0 / 3.3, and the noise bias went from +0.34 to -0.005 in log10. That is the paragraph
% below, and it answers the open question nomenclature.md left on whether INR separates from NR.
% 0ffbda7, NMR from 87889e6, the rest from 1c2ae6f); DATA_DIRS is a search path, so adding 87889e6
% was the whole change. The IR correlation corner that used to be panel B is out of the body figure
% and is a supplement. Panel size in the .Rmd is measured, not set: it is what the seven columns
% leave once the 7 pt type has taken its space, so raising the type re-sizes the panels by itself.
% The file is 7.00 x 6.07 in and renders at 7.425 x 6.44 in, since width=\linewidth inside fullwidth
% enlarges it by 6.1%.
\includegraphics[width=\linewidth]{Figure_2}
\end{fullwidth}
\caption{\textbf{Parameter recovery and the calibration of the reported interval, for the whole ladder at one design cell.} $N_{\mathrm{ch}}=100$, noise label $0.1$ ($\tilde S = 0.01$), $\Delta\!\cdot\!k_{\mathrm{off}}=0.1$; $10{,}000$ recordings in groups of ten, so one thousand maximum-likelihood estimates per member. Columns are the seven rungs in ladder order, from least squares to the boundary-conditioned filter. Rows are three parameter pairs, each chosen because a different property separates the members in it: the kinetic pair ($\log_{10}k_{\mathrm{on}}$, $\log_{10}k_{\mathrm{off}}$), which the deterministic transient determines and on which every member including least squares competes; the amplitude pair ($\log_{10}N_{\mathrm{ch}}$, $\log_{10}i$), which only the gating fluctuations separate; and the noise pair ($\log_{10}\sigma$, $\log_{10}N_{\mathrm{ch}}$), where $\sigma$ is the instrumental noise level, a free parameter of every likelihood member, with a simulated value of $10^{-4}$. Each panel carries the cloud of estimates and three $95\%$ ellipses drawn at the cloud mean: the empirical covariance, the Fisher covariance the member reports, and the sandwich covariance. The two annotated numbers are the area of the empirical ellipse over the area of the reported one (orange) and over the corrected one (blue), so a member is calibrated when both read $1.0$. Markers are the simulation truth, the mean of the cloud, and the truth displaced by the predicted bias. The two least-squares cells that are empty are empty by construction and say so: that configuration holds the unitary current and the noise level fixed at their simulated values, so its estimates would lie on a line and the area ratio would be $0/0$. The third row separates the members in a way the other two do not. The open-loop member places its noise estimate at $\log_{10}\sigma = -3.67$ against a truth of $-4$, a factor of $2.2$ too much instrumental noise; every other member sits within $0.04$ of the truth, the one-endpoint non-recursive member included. % src: projects/eLife_2025/figures/paper_both/figure_2.Rmd (mean log10 Current_Noise: NR -3.667, NMR -3.662, R -4.023, MR -4.036, VR -4.031, IR -4.005)
Covariances are the Gaussian-Fisher family at the pooled estimate, scaled by the reciprocal of the group size.}
\label{fig:clouds}
\end{figure}

\subsection*{The calibration cascade in time}

The miscalibration of Figure~\ref{fig:clouds} is a property of the likelihood and not of the optimiser, and Figure~\ref{fig:time} shows where in the recording it is produced. One common set of a thousand recordings at the same cell ($N_{\mathrm{ch}}=100$, noise label $0.1$, $\Delta\!\cdot\!k_{\mathrm{off}} = 0.1$) is scored interval by interval by each member, with no fitting involved. The least-squares arm here is the six-parameter configuration, which keeps the unitary current and the noise level free so that its per-interval quantities align index for index with the macroscopic dumps (Methods); nothing measured in this subsection may be carried across to the four-parameter configuration of the later figures.

The members separate first on the log-likelihood they assign to the same data. The ensemble-mean total is $-262 \pm 0.73$ for \texttt{LSE}, $-230 \pm 0.83$ for \texttt{NR}, $-147 \pm 0.20$ for \texttt{R} and $-143 \pm 0.24$ for \texttt{IR}, a spread of $118.9$ nats from the bottom of the ladder to the top. % src: projects/eLife_2025/figures/paper_both/figure_3.html (caption-numbers chunk of figure_3.Rmd; 1000 recordings)
The ordering follows the ladder, and most of the span is bought by modelling the gating fluctuations at all.

The hinge of the argument is that per interval every member reports its information correctly. The variance of the per-interval score matches the per-interval Fisher information for all four, and the median of $\log_{10}(J_t/F_t)$ over the informative steps is $+0.184$ for \texttt{LSE}, $-0.038$ for \texttt{NR}, $-0.181$ for \texttt{R} and $+0.006$ for \texttt{IR} on $k_{\mathrm{on}}$, with $-0.145$ and $+0.033$ for \texttt{R} and \texttt{IR} on $k_{\mathrm{off}}$. % src: projects/eLife_2025/figures/paper_both/figure_3_supplement_1.html (caption-numbers chunk, D-G table; median per-step log10(J_t/F_t) over the informative steps, all four parameters. Re-measured from the digests 2026-07-31 and reproducing the archived notebook to five decimals: k_on 0.183889 / -0.037663 / -0.181439 / +0.005922, k_off -0.144531 / +0.033493)
Accumulated over the recording the same ratio reads $14.96$ ($95\%$ interval $13.86$ to $15.93$) for \texttt{LSE}, $13.63$ ($12.52$ to $14.95$) for \texttt{NR}, $1.07$ ($0.975$ to $1.17$) for \texttt{R} and $1.08$ ($0.995$ to $1.17$) for \texttt{IR}. % src: projects/eLife_2025/figures/paper_both/figure_3.html (caption-numbers chunk; accumulated J_T/F_T with bootstrap interval)
The two non-recursive members are within a few percent of the identity at each interval and off by more than an order of magnitude once the intervals are summed. The open-loop member also changes the sign of the discrepancy between the per-step and the accumulated statistic, and a per-sample modelling error cannot do that. Least squares does not change sign, sitting above the calibrated value at both levels, which is the signature of its constant-variance premise rather than of accumulation. % src: projects/eLife_2025/figures/paper_both/figure_3_supplement_1.html (per-step median log10(J_t/F_t): LSE +0.184, NR -0.038, R -0.181, IR +0.006). The constant-variance reading is now SHOWN and not only asserted: row E of that supplement is one curve and not four in the least-squares column, because with a single variance the parameter cancels out of Var(s_t)/F_t (median 0.183889 for all four, identical to six decimals)

The least-squares ratio needs the qualification the Diagnostics section attaches to it. The Gaussian Fisher information of a least-squares fit presumes one constant residual variance, and the exact simulator violates that premise by construction, because the variance of a recorded interval rises and falls with the open population. The factor of fifteen is the size of that break, and it is not an information-matrix-equality failure of the same kind as the open-loop gating likelihood's, where the predictive variance does follow the gating. The two ratios are reported side by side and are not read as two values of one comparable quantity.

The same checks run for the parameters the figure has no colours for, and for the first moment of the residual, in Figure~\ref{fig:time}--figure supplement 2. Two things there are worth carrying back: the instrumental noise level is the only parameter measured before the agonist arrives, and its score is white even in the members whose kinetic directions are not, so the memory that carries the accumulated error is a property of the gating directions; and least squares has no score in that direction at all, because its noise level is a plug-in from the same record rather than a free parameter. % src: projects/eLife_2025/figures/paper_both/figure_3_supplement_3.html (caption-numbers chunk: noise informative at 100 of 100 intervals with 20 pre-agonist against 0 for k_on and i; lag-one noise score ACF -0.004 NR, -0.008 R, both covering zero; LSE noise informative at 0 intervals)

The mechanism is visible in the same figure. At lag one the autocorrelation of the per-interval score on $k_{\mathrm{off}}$ is $-0.0043 \pm 0.0043$ for \texttt{IR}, indistinguishable from zero, against $0.191 \pm 0.004$ for \texttt{R}, $0.864 \pm 0.0019$ for \texttt{NR} and $0.856 \pm 0.0022$ for \texttt{LSE}. % src: projects/eLife_2025/figures/paper_both/figure_3.html (caption-numbers chunk; lag-one score ACF, k_off)
A positively correlated score makes the variance of its sum exceed the sum of the per-interval informations, which is exactly how the per-interval identity can hold while the accumulated one fails. The temporal correlation the two non-recursive members leave in their score is the local time correlation of the current that \citet{milescu2005maximum} named as the dominant error of independent-interval fitting, measured here at the level of the inference rather than the data.

Row C of Figure~\ref{fig:time} carries a second result, and it belongs to the measurement rather than to any member. The information about the channel number falls to zero the moment the agonist is removed and the open population relaxes: having conditioned on everything up to that point, the filter gains nothing further from watching the channels decay. The closing rate is the exception and stays informative through the relaxation, which its own kinetics govern. The opening rate and the unitary current behave like the channel number (Figure~\ref{fig:time}--figure supplement 1).
% TODO(Luciano) 2026-07-31 -- MEASURED, and the unitary current does NOT behave like the channel
% number. Mean F_t over the last ten intervals, IR: k_on 7.9e-44, N_ch 1.4e-21, unitary current 0.66,
% k_off 1.13. So k_on and N_ch reach the numerical floor within six intervals of the washout and the
% unitary current does not: it falls by a factor of 36 and levels at the same order as k_off. The two
% survivals are of different kinds, which is the better sentence: 81% of what k_off keeps comes from
% the MEAN (the shape of the decay it governs) and 71% of what the unitary current keeps comes from
% the VARIANCE (the gating fluctuation of the channels still open), while 99.6% of N_ch's is from the
% mean and it is that mean derivative that vanishes once the filter has conditioned on the past.
% Second point, same sentence in Discussion: "the shape is the same in every column" holds for the
% conditioned members only. LSE and NR retain 15-16% of their pulse information about N_ch through
% the washout against 0.5-0.8% for R and IR, which is what this paragraph's own mechanism predicts,
% since gaining nothing further from the decay requires having conditioned on it.
% src: projects/eLife_2025/figures/paper_both/figure_3_supplement_1.html (washout table, row C).
% NOT EDITED: the correction changes a stated result, so it is yours to make. The shape is the same in every column that carries the comparison, which is what makes it a property of the macroscopic observable rather than of one filter, and it answers directly the question of when each parameter stops being measured that a recent treatment of this class of models leaves open \cite{delcore2025parameter}.
% NOTE (LM): the generality claim is now testable in the body, since row C runs across all four columns. Verify it against the render before submission: if the non-recursive columns depart, the result is a property of the recursive filters and this sentence has to say so. The mechanism given here (having conditioned on everything up to that point) is a recursion argument, so a departure would not be surprising. Least squares carries no channel-number curve, so that column cannot support this claim either way.

\begin{figure}[htbp]
\begin{fullwidth}
\centering
\includegraphics[width=\linewidth]{Figure_3}
\end{fullwidth}
\caption{\textbf{The calibration cascade in time, over the whole implemented family.} One set of $1{,}000$ recordings ($N_{\mathrm{ch}}=100$, noise label $0.1$, $\Delta\!\cdot\!k_{\mathrm{off}}=0.1$) is scored by each of the seven members (columns, in cost-ladder order), with no fitting. Rows A and B are properties of the data, in grey; rows C to G are per parameter, the closing rate and the channel number in colour. (A) An example recording with the predicted current, its band the filter's own predictive $\pm1$ s.d., annotated with the ensemble-mean log-likelihood and its standard error. (B) The mean squared standardized residual, one when the predictive variance is right. (C) The per-interval Fisher information, so where each parameter is measured. (D) The mean per-interval score over $\sqrt{F_t}$, zero when unbiased and standardized so parameters with different units share an axis. (E) The per-interval information ratio $\mathrm{Var}(s_t)/F_t$, one at each interval. (F) The same ratio accumulated, $J_T/F_T$. (G) The autocorrelation against lag of the standardized residual, in black, and of the score for each parameter, in its colour; the score's is why (E) holds where (F) does not, and the residual's is the same memory in the data and the only quantity here an experimentalist can compute on a real recording. Read across the columns the figure is a factorial on the two switches of the lattice. Interval averaging alone (\texttt{INR}) repairs the data level and not the accumulation, $0.9995$ on (B) against $20.8$ on (F); recursion alone (\texttt{R}) repairs the accumulation and not the data level, $0.782$ against $1.07$; only the two together repair both. The one-endpoint members are the other half of that reading: \texttt{MR} is worse than \texttt{R} rather than between it and \texttt{IR}, so conditioning on a single interval end does not interpolate, and \texttt{VR}, which removes the between-endpoint variance without restoring it in the gain, is worse still. Bands are $95\%$ bootstrap intervals over the recordings; curves stop where the information reaches zero. Non-negative quantities are on a logarithmic axis with the calibrated value on the ticks, signed quantities on a linear axis with zero on the ticks; the three open-loop columns share a y-scale block distinct from the four recursive ones. The channel number is not drawn for least squares, which cannot separate it from the unitary current.}  % src: projects/eLife_2025/figures/paper_both/figure_3.Rmd (promoted 2026-07-31 from the four-column version, now archived; seven columns, seven rows, two parameters) ; Figure_3_caption.md (the long form this compresses)
\label{fig:time}
\figsupp{\textbf{Per-step information against score variance, for every member and all four identified parameters.} All seven members across the columns and the four parameters down the page as bands of two rows: $k_{\mathrm{on}}$ (A, B), $k_{\mathrm{off}}$ (C, D), the unitary current (E, F) and $N_{\mathrm{ch}}$ (G, H). The upper row of each band carries the per-interval Fisher information and the score variance as magnitudes, which lie on each other where the member is calibrated at that interval; the lower row is $\log_{10}(J_t/F_t)$ with its $95\%$ bootstrap interval, zero when they do. Per-interval calibration is not monotone along the cost ladder: on $k_{\mathrm{off}}$ the median reads $+0.184$ for \texttt{LSE}, $-0.038$ for \texttt{NR}, $+0.007$ for \texttt{INR}, $-0.145$ for \texttt{R}, $-0.229$ for \texttt{MR}, $-0.135$ for \texttt{VR} and $+0.033$ for \texttt{IR}, so the cheapest member that models the interval average is better calibrated interval by interval than any of the three recursive members that cost more. Curves stop at the information mask, which is also where the upper axis ends. In the two recursive columns the $k_{\mathrm{on}}$ and $N_{\mathrm{ch}}$ bands reach that floor within six intervals of the agonist being removed while the $k_{\mathrm{off}}$ and unitary-current bands do not. The least-squares ratio row is the same curve in all four bands, with a median of $0.183889$ in every one: with a single constant variance the parameter cancels out of the ratio, which leaves the true residual variance over the assumed one. The channel-number band carries no least-squares panel, which with the gating fluctuations unmodelled cannot separate that parameter from the unitary current.}{\includegraphics[width=\linewidth]{Figure_3_supplement_1}}
\figsupp{\textbf{The checks this figure has no room for.} All seven members at the same cell, carrying (A) the mean standardized residual, the one of the three properties of $r_t$ not drawn above; (B) the score bias for the three parameters not coloured above, $k_{\mathrm{on}}$, the unitary current and the instrumental noise level; and (C) their score autocorrelation. The first residual moment separates nobody, which is what makes the other two properties the informative ones. In (B) the ordering repeats the one in figure supplement 2: counting informative intervals whose interval excludes zero against a chance level of $0.05$, \texttt{IR} reads $0.067$, $0.050$ and $0.050$, \texttt{INR} $0.100$, $0.100$ and $0.000$, and the three recursive members between $0.409$ and $0.975$. The noise level is the control direction: it is the only parameter measured before the agonist arrives, and its score is white even in the members whose kinetic directions are not, so the memory that carries the accumulated information error lives in the gating directions and not in the instrumental one. Least squares carries no noise curve because its noise level is a plug-in from the same record rather than a free direction, so its score there vanishes identically, which is the same fact as its mean squared residual being pinned at one.}{\includegraphics[width=\linewidth]{Figure_3_supplement_2}}
% SUPPLEMENT ORDER, set 2026-07-31: declaration order must match first-citation order in the text.
% The residual triptych is cited at 03_diagnostics.tex:7, the four-parameter widening at
% 03_diagnostics.tex:47 and twice more; Diagnostics precedes this section. The family widening is
% cited nowhere and takes the last slot.
\end{figure}

\subsection*{The design plane, in both moments}

One cell is not a result. Figure~\ref{fig:plane} repeats the two moments over the whole grid, with the members across the columns, the channel count down the rows, and the acquisition interval against the instrumental noise inside each panel. It splits the display by moment rather than by member, so the least-squares arm stays alongside the likelihoods it is being compared with. Half (A) is evaluated at the simulation truth, where a bias is visible; half (B) at the pooled optimum, where the score vanishes by construction and the information comparison is clean (Methods). Both halves are drawn on one scale, the factor by which the report is wrong, so the paleness of (A) against the saturation of (B) is a comparison and not an impression. The least-squares arm here is the four-parameter configuration, which holds the unitary current and the instrumental noise level fixed at the values the recordings were simulated from (Methods), so it holds one column per half, the closing rate only: with the unitary current assumed, the amplitude ridge that limits the other three members does not exist for it, and a channel-number panel would read as least squares winning a contest it did not enter.

More channels do not repair the recursive instantaneous filter.
% TODO-SIGN: direction convention as in Diagnostics.
Taking the median over the seven intervals at noise label $0.1$, the $k_{\mathrm{off}}$ distortion of \texttt{R} sits on a flat floor of $1.372$, $1.440$, $1.439$ and $1.42$ across $N_{\mathrm{ch}}$ $10$, $100$, $1000$ and $10^4$, while \texttt{IR} converges toward one over the same range, $1.324$, $1.099$, $1.000$, $1.00$. % src: projects/eLife_2025/figures/paper_both/figure_4_distortion.html (caption-numbers chunk; median over 7 intervals, noise 0.1)
The open-loop member moves the other way, $21.92$, $16.83$, $31.32$ and $76.52$, and the least-squares arm holds a floor an order of magnitude above one, $11.79$, $13.06$, $13.04$ and $12.97$. % src: projects/eLife_2025/figures/paper_both/figure_4_distortion.html
In the channel-number direction the two recursive members separate in sign as well as in size: \texttt{IR} runs $0.851$, $0.936$, $0.987$, $1.00$ and stays on the conservative side, while \texttt{R} crosses from conservative to over-confident, $0.866$, $1.100$, $1.230$, $1.26$, and keeps rising. % src: projects/eLife_2025/figures/paper_both/figure_4_distortion.html
The two members agree at the fewest channels and separate from there, so the $10^4$ column is the informative one.

The first moment carries a separate cost. The recursive instantaneous filter has a bias in the channel number of $0.126$, $0.148$, $0.138$ and $0.138$ in $\log_{10}$ units across the same four decades, close to a $35$ to $40\%$ error that does not shrink as channels are added; the open-loop member carries $0.073$ to $0.084$ in the same direction, and the boundary-conditioned member is at or near zero everywhere. % src: projects/eLife_2025/figures/paper_both/figure_4_bias.html (caption-numbers chunk; median over 7 intervals, noise 0.1)
Least squares has a rate bias of essentially zero over the same plane. % src: projects/eLife_2025/figures/paper_both/figure_4_bias.html (LSE k_off bias 0.00891 / 0.000164 / 0 / 0)
That concession is worth making plainly: the classical fit recovers the rates, and what it gets wrong is how well it says it recovered them.

The two recursive members also differ in where on the plane their departures sit, and that matters as much as their size. Counting the fraction of cells that depart from one by at least $15\%$, parameter by parameter, \texttt{R} departs on $0.4375$, $0.5625$, $0.2232$, $0.3125$ and $0.0000$ of the plane for $k_{\mathrm{on}}$, $k_{\mathrm{off}}$, $i$, $N_{\mathrm{ch}}$ and the noise level, against $0.0179$, $0.0804$, $0.0000$, $0.0268$ and $0.0000$ for \texttt{IR}. % src: projects/eLife_2025/figures/paper_both/figure_4_supplement_2.html (fraction of cells departing >= 15%, per parameter and algorithm)
The distortion of the instantaneous filter is spread across parameters and across the plane; the distortion of the boundary-conditioned filter is confined to the few-channel corner. Conditioning localises the error as well as shrinking it.

Instrumental noise moves the whole picture continuously, and it moves the window-ignoring members furthest. At $N_{\mathrm{ch}}=10$ the $k_{\mathrm{off}}$ distortion of least squares runs $11.79$ at noise label $0.1$, $1.14$ at label $100$ and $1.00$ at label $1000$; at $N_{\mathrm{ch}}=10^4$ the same sequence runs $12.97$, then $2.54$ at label $10^4$, $1.15$ at $10^5$ and $1.00$ at $10^6$. % src: projects/eLife_2025/figures/paper_both/figure_4_distortion.html ("DOES NOISE CURE IT" table, k_off medians by noise and N_ch)
The recursive instantaneous filter follows the same pattern about three decades lower, running $1.42$, $1.03$ at label $100$ and $1.00$ at label $1000$ at $N_{\mathrm{ch}}=10^4$. % src: projects/eLife_2025/figures/paper_both/figure_4_distortion.html
The reason is that instrumental noise adds a variance that carries no gating memory, so as it grows it dilutes the temporal structure a window-ignoring likelihood fails to model, until there is nothing left to fail on. The channel count at which each member's error bar becomes honest therefore rises with the noise level, which is the crossover Figure~\ref{fig:usage} draws as a boundary.

Over the whole measured grid, counting every parameter at every cell, the reported uncertainty is within $\pm15\%$ of the delivered one for $94\%$ of the boundary-conditioned member's points ($1977$ of $2100$), $70\%$ of the recursive instantaneous member's ($625$ of $896$), and $33\%$ of the open-loop member's ($277$ of $840$); the one-endpoint member carried in the supplement reaches $20\%$. Least squares carries no entry in this comparison, because its distortion matrix has four dimensions against the likelihood members' six and the two counts are not commensurable (Methods). % src: papers/1_method/decisions/recompute/d4_distortion_envelopes.py re-run over all three data directories 1c2ae6f + 87889e6 + 0ffbda7 (the committed copy of the script lists only the first two, which is why D-4 §1 still reports 93 / 51 / 20 / 0)
The coverage is honestly unequal and has to be quoted with the number: the boundary-conditioned member spans eleven channel counts from $5$ to $10^4$ across eleven noise levels, not on a full rectangle, and the others cover less. % src: projects/eLife_2025/figures/data/{1c2ae6f,87889e6,0ffbda7} (macro_IR nsim_10000 cells) The open-loop member's $33\%$ is worth stating explicitly, because it recovers across the high-noise crossover and is therefore calibrated somewhere, which a flat statement that it never is would get wrong.

The boundary-conditioned member is itself an approximation and fails in both directions. Over the whole measured grid its distortion runs from $0.645$, at $N_{\mathrm{ch}}=10$, noise label $0.05$, $\Delta\!\cdot\!k_{\mathrm{off}} = 0.01$, on the channel-number direction, where it over-states its own uncertainty, to $1.701$, at $N_{\mathrm{ch}}=5$, noise label $0.1$, $\Delta\!\cdot\!k_{\mathrm{off}} = 1$, on the closing rate, where it under-states it. % src: papers/1_method/decisions/recompute/d4_distortion_envelopes.py (min and max of the diagonal distortion over the three data directories)
Both excursions are in the few-channel corner, which is where the occupancy closure of the Theory section is expected to degrade.

The mechanism behind the flat floor is the split of the distortion into a per-sample part and a correlation part (Figure~\ref{fig:plane}--figure supplement 3). On the closing rate at noise label $0.1$, the per-sample part goes to one as channels are added for both recursive members, $1.09$ to $0.998$ for \texttt{R} and $1.10$ to $1.007$ for \texttt{IR} across the four decades, so the single-interval Gaussian approximation becomes exact for both at the same rate. % src: projects/eLife_2025/figures/paper_both/Figure_4_supplement_3_caption.md (caption-numbers chunk of figure_4_supplement_3.Rmd; sample-part medians)
What separates them is the correlation part, which rises for \texttt{R} from $1.35$ to $1.47$ and falls for \texttt{IR} from $1.24$ to $1.00$. % src: projects/eLife_2025/figures/paper_both/Figure_4_supplement_3_caption.md (correlation-part medians)
The two members are equally faithful to a single interval, and they part company on how much temporal dependence they leave in the score, which is what conditioning on both interval endpoints removes.

The same memory is visible in the data, without any parameter. The integrated autocorrelation of the standardized residual, which counts how many intervals a likelihood effectively has for every interval it believes it has, reads $1.008$ for \texttt{IR}, $1.30$ for \texttt{R}, $3.93$ for \texttt{NR} and $9.26$ for \texttt{LSE} at noise label $0.1$. % src: projects/eLife_2025/figures/paper_both/figure_4_supplement_8.html (residual integral-lag by algorithm and noise, median over N_ch and interval)
The classical fit therefore behaves as though it had about nine times more independent observations than the recording supplies, which is the degrees-of-freedom cost of fitting correlated data as though it were independent.

Two features of the plane are easy to misread and belong in the text. The first is the left edge. Where the acquisition interval is far shorter than the channel's correlation time, successive samples carry redundant information, and a faithful likelihood has to register that redundancy, so a large correlation distortion at the shortest intervals is the correct cost of oversampling rather than a defect. The second is a single cell carried openly. At $N_{\mathrm{ch}}=10^4$ and $\Delta\!\cdot\!k_{\mathrm{off}}=1$ the amplitude directions of the recursive instantaneous filter are not identified: the Gaussian Fisher information is ill-conditioned there, so a displacement that costs almost no likelihood is amplified into a large parameter shift, and the reported bias runs off the scale of the map. Those cells are drawn grey. Their sign and their existence are real; their magnitude is a property of the conditioning and is not a physical bias. % src: projects/eLife_2025/figures/paper_both/figure_4_supplement_1.html (greyed cells: 8 -> R unitary_current and R Num_ch_mean at N_ch 10000) ; figure_4_common.R:212-216 (grey = off-scale AND Fisher condition number above 3e4)

\begin{figure}[htbp]
\begin{fullwidth}
\centering
% One file since 2026-07-30 (figure_4.Rmd, which replaces the stacked byNch pair): the moment is
% the outermost column grouping rather than a second file, which halves the height and puts both
% moments on one shared scale. 7.4 x 7.3 in, isotropic in the two axes.
\includegraphics[width=\linewidth]{Figure_4}
\end{fullwidth}
\caption{\textbf{The design plane: how wrong each member's report is, in both moments, on one scale.} Ten thousand recordings are simulated at every cell of the grid and scored by each member of the cost ladder. (A) The distortion-induced bias, evaluated at the simulation truth $\theta_{\mathrm{sim}}$: how far the estimate is displaced. (B) The information distortion, evaluated at the pooled optimum $\theta_{\mathrm{pool}}$, where the score vanishes by construction so a displaced gradient cannot contaminate it: how far the uncertainty a member reports is from the one it delivers. Columns nest twice, member and then parameter, $k_{\mathrm{off}}$ and $N_{\mathrm{ch}}$. Least squares holds one column per half, because its four-parameter configuration fixes the unitary current and the noise level at the simulated values (Methods), so a channel-number comparison against it would not be like for like. Rows are the channel count, each given a height proportional to the range of instrumental noise the sweep reaches there, so a decade of noise is the same distance in every row. Inside a panel the axes are the sampling interval in units of $\tau$ and the dimensionless instrumental noise, and one decade is the same physical distance on both, so a field's shape can be read as a shape. Colour is shared by the two halves: the factor by which the report is wrong, centred on one, pale right and saturated wrong in either direction, blue conservative. What a factor means differs between them and the difference is not cosmetic: in (A) it is a factor on the parameter, so $2$ means the estimate is out by a factor of two, while in (B) it is a ratio of variances, so $2$ means the reported error bar is $\sqrt{2}$, about $40\%$, too narrow. The scale is clipped at $10^{\pm3}$, which touches $22$ of $3276$ cells, twenty of them the open-loop member, where being wrong by $10^3$ and by $10^{12}$ say the same thing. Lines carry two codes: colour says which quantity, dark hairlines cutting the colour field and white lines on a dark casing cutting the distortion-corrected standard error, which asks whether the parameter can be measured at all rather than whether its error bar is honest; dash says which criterion, solid a factor $1.15$ and dashed a factor $2$, the two thresholds used throughout. White rather than a second hue because the field is red and blue, so hue would fail for red-green colour blindness where the lines matter most. Both criteria are edges of the colour bands, which is why the dark pair needs no key of its own.} % src: projects/eLife_2025/figures/paper_both/Figure_4_caption.md (the long form this compresses) ; figure_4.Rmd (layout, shared factor scale, the two line codes, the clipping count) ; figure_4_common.R (EXCLUDE_ROWS, nch_noise_span, CRIT_SOL/CRIT_DSH)
\label{fig:plane}
\figsupp{The distortion-induced bias for all five parameters, the recursive instantaneous filter against the boundary-conditioned one: the departures of the first are spread across parameters, those of the second are confined.}{\includegraphics[width=\linewidth]{Figure_4_supplement_1}}
\figsupp{The information distortion for all five parameters, the same pair, showing the same contrast in the second moment.}{\includegraphics[width=\linewidth]{Figure_4_supplement_2}}
% NOTE: the map version of the sample/correlation split (figure_4_supplement_3.Rmd) is not built;
% the line version below carries the same decomposition and occupies the third supplement slot, which
% is the slot the Diagnostics section cross-references for the split.
\figsupp{The distortion decomposed into a per-sample part and a correlation part, against channel count and resolved by instrumental noise. The per-sample parts of the two recursive members coincide and go to one as channels are added, while the correlation parts separate.}{\includegraphics[width=\linewidth]{Figure_4_supplement_4}}
\figsupp{The sandwich prediction against the empirical multivariate distribution of the estimates, as a Mahalanobis quantile-quantile plot with a $\chi^2$ reference. The reported Fisher covariance under-covers the nominal $95\%$ region and the sandwich covers it.}{\includegraphics[width=\linewidth]{Figure_4_supplement_5}}
\figsupp{How near-singular the Gaussian Fisher information is over the same plane, which is what decides where a quantity is left undefined rather than set to zero.}{\includegraphics[width=\linewidth]{Figure_4_supplement_6}}
\figsupp{The distortion-corrected standard error over the same plane: what each member achieves, as opposed to what it reports.}{\includegraphics[width=\linewidth]{Figure_4_supplement_7}}
\figsupp{The integrated autocorrelation of the standardized residual over the plane: the gating memory that least squares and the open-loop member leave in the data.}{\includegraphics[width=\linewidth]{Figure_4_supplement_8}}
\figsupp{The integrated autocorrelation of the score over the plane, for $k_{\mathrm{off}}$ and $N_{\mathrm{ch}}$: the same memory at the level of the inference, which is what carries the accumulated information error.}{\includegraphics[width=\linewidth]{Figure_4_supplement_9}}
\figsupp{The whole-matrix distortion, as the affine-invariant distance of the information-distortion matrix to the identity. The ordering across members is the one the per-parameter maps show; the scalar itself accumulates over eigenvalues and the least-squares matrix is $4\times4$ against the family's $6\times6$, so the values are read as a ranking and not as a common scale.}{\includegraphics[width=\linewidth]{Figure_4_supplement_10}}
\figsupp{The same whole-matrix distance, separated into its per-sample and correlation causes over the design space, with the same dimension caveat.}{\includegraphics[width=\linewidth]{Figure_4_supplement_11}}
\end{figure}

\subsection*{The information budget per parameter}

A calibrated member's reported covariance can be spent, and Figure~\ref{fig:design} spends it. Only the boundary-conditioned member is shown, because the quantity is a covariance the member reports and only a member whose report has been checked can guide a design. The covariance used is the distortion-corrected one, so what is plotted is the precision the experiment achieves.

The quantity is a trade-off at a fixed channel budget. Plotting $N_{\mathrm{ch}}$ times the corrected variance of a parameter gives, up to the fixed budget constant, the variance of the estimate pooled over as many recordings as the budget allows. A flat line means the information is linear in channels, so it makes no difference whether the budget is spent as a few large patches or many small ones. A rising line means concentrating the budget into few large recordings yields less information.

The rates and the amplitude directions give opposite advice. Over the four decades of channel count at noise label $0.1$, the pooled variance grows by a factor of $1.24$ for the closing rate and $3.11$ for the opening rate, against $288.19$ for the unitary current, $60.50$ for the channel number and $3236.85$ for the instrumental noise level. % src: projects/eLife_2025/figures/paper_both/figure_5.html (caption-numbers chunk; fan factor, N_ch*corrected covariance at N_ch 10^4 over N_ch 10, noise 0.1)
For the rates, a channel budget can be concentrated at almost no cost. For the amplitude, the count and the noise level, the same budget spent as many small recordings gives a substantially better estimate. A noisier rig penalises the few-channel end, where the instrumental noise is large next to the gating variance, and it does so first in the rates, which are otherwise insensitive to how the budget is split. The acquisition interval runs the same way in every panel, with the pooled variance smallest at short intervals and growing toward the coarse end, mildly for the rates and more for the amplitude directions, which agrees with the calibration trade-off of Figure~\ref{fig:plane}.

The scope of this figure is narrower than the trade-off it states. Everything here is the non-stationary protocol, a single concentration jump with a rising and a relaxing phase. The rates are recoverable from the shape of that deterministic transient, which is why they cost so little. Under a stationary protocol the mean current is flat, the transient carries no rate information, and the balance between the rates and the amplitude directions inverts. The advice below should not be carried across that boundary.

\begin{figure}[htbp]
\begin{fullwidth}
\centering
\includegraphics[width=\linewidth]{Figure_5}
\end{fullwidth}
\caption{\textbf{Where the boundary-conditioned filter stops reporting an honest error bar, and which of the two discarded moments is responsible.} The same measurement as Figure~\ref{fig:plane}, the Gaussian information distortion at the pooled optimum on the same runs, resolved on the corner where \texttt{IR} departs: on Figure~\ref{fig:plane}'s scale a factor of $1.5$ is a faint tint and here it needs an axis. The three blocks are the distortion and the two parts it decomposes into: (A) the total, (B) the part carried by the per-sample non-Gaussianity of the score, and (C) the part carried by its correlation across intervals; their product reproduces the total to within $1\%$ over the whole grid. Each block is a grid of the design plane in the orientation of Figures~\ref{fig:plane} and~\ref{fig:usage}: the channel count increases to the right, the dimensionless noise $\tilde S$ increases upward, and inside every cell the horizontal axis is $\Delta\!\cdot\!k_{\mathrm{off}}$ from $0.01$ to $1$, shared by every panel and labelled once per block. The vertical axis is the distortion factor on a logarithmic scale, shared by all twelve columns so that they can be compared, with the green line at $1$ marking a report that is neither over- nor under-confident. Colour is the parameter. Bands are $95\%$ percentile bootstrap intervals over recordings, of median half-width $0.03$. \textbf{The departure is two-sided within one cell:} at ten channels and $\tilde S = 0.005$ the reported error bar of $k_{\mathrm{off}}$ is too small by a factor of $1.62$ at the shortest interval, while that of $N_{\mathrm{ch}}$ is too large by a factor of $0.65$, so the same likelihood is over-confident about the rate and conservative about the count at once. \textbf{The two parts answer to the interval in opposite directions:} across that cell the correlation part of $k_{\mathrm{off}}$ falls from $1.52$ to $1.27$ as the interval coarsens while the per-sample part rises from $1.06$ to $1.19$, so sampling coarsely buys back the one and not the other. Both parts vanish by $\tilde S = 1$, the instrumental noise making the emission more Gaussian. Two limits: the family applies both of its Gaussian approximations at every step and the score sees only their product, so a departure measured here cannot be assigned to one of them; and the grid is truncated where the effect grows, the departure of $k_{\mathrm{off}}$ still accelerating at the lowest simulated noise ($1.22$, $1.37$, $1.52$ over the last three noise steps at ten channels), so the figure locates the onset and not the far side.} % src: projects/eLife_2025/figures/paper_both/Figure_5_caption.md (full caption, with the component paths and the composition check) ; figure_5.Rmd
% NOTE: the label is still fig:design, from the information-budget figure this float used to hold.
% Renaming it to fig:limit needs the "information budget" subsection above rewritten, since its
% prose introduces the demoted figure; do both in one pass.
\label{fig:design}
\end{figure}

\subsection*{A usage map for a macroscopic recording}

Figure~\ref{fig:usage} collects the measurements into one plane, the channel number against the instrumental noise, at a single acquisition interval $\Delta\!\cdot\!k_{\mathrm{off}}=0.1$. Unlike the earlier maps, its content is the boundaries: the measured cells are the input and the lines are the output. Four boundaries are drawn, and they are two questions asked of two methods. Two ask whether a method's reported error bar is honest, one for the boundary-conditioned likelihood and one for least squares, and are drawn where that method's information distortion crosses $1.15$, a $7\%$ error on the reported standard deviation. Two ask whether a parameter is reachable at all, one for the amplitude pair and one for the rates, and are drawn where the distortion-corrected standard error crosses a factor of two. The rate boundaries are taken from the least-squares arm and the amplitude boundaries from the likelihood, because the least-squares arm of this figure is the four-parameter configuration, which holds the unitary current fixed and so does not estimate it (Methods), while the two arms agree on the rates over the channel counts where both are measured and least squares is the arm with the high-noise coverage. % src: projects/eLife_2025/figures/paper_both/figure_6.Rmd:31 (rate boundaries sourced from LSE for coverage; amplitude pair from IR)

Read upward at a fixed channel count, the recording loses capabilities in a fixed order, which gives five regions. In region~0, below the lowest boundary, the Gaussian closure is itself misspecified, so the likelihood's own error bar is wrong and the recording is not a macroscopic problem. Region~1 delivers the unitary current, the channel number and an honest interval, while least squares reports an interval several times too narrow. In region~2 the amplitude pair is no longer separable and the rates alone survive, so the likelihood still buys an honest interval. In region~3 least squares is calibrated too and buys the same answer more cheaply. Above region~3 nothing useful is left. The boundary of region~0 is the one with a negative slope, measured at $-0.56$ in $\mathrm{d}\log_{10}\mathrm{noise}/\mathrm{d}\log_{10}N_{\mathrm{ch}}$, because instrumental noise makes the emission more Gaussian and so relaxes the occupancy closure while hurting everything else. % src: projects/eLife_2025/figures/paper_both/figure_6.html (fitted log-log slopes at Delta.k_off = 0.1)

The boundary at which the classical error bar becomes honest is the one this paper predicts and then measures. In noise-label units it sits at $164$ and $130$ at $N_{\mathrm{ch}}=10$, at $1.28\times10^{3}$ and $1.53\times10^{3}$ at $100$, at $1.54\times10^{4}$ and $1.51\times10^{4}$ at $1000$, and at $1.62\times10^{5}$ and $1.59\times10^{5}$ at $10^4$, taking the channel-number and the closing-rate directions in that order; the dimensionless noise is one tenth of each. % src: projects/eLife_2025/figures/paper_both/figure_6.html (crossing table, q = lse, Delta.k_off = 0.1)
Fitted in logarithms, those crossings give slopes of $1.01$ and $1.03$, so the boundary runs as noise proportional to channel number. % src: projects/eLife_2025/figures/paper_both/figure_6.html (fitted log-log slopes)
That is the crossover the mechanism predicts, since the gating variance grows with the channel count while the instrumental variance does not, so the ratio of the two is what decides whether a window-ignoring fit is honest. Here it is a measured slope rather than an assumed scaling.

The ordering of the boundaries is the result the map exists to state. There is no region of the measured plane in which the amplitude still carries information and the classical error bar is simultaneously trustworthy; the two boundaries are separated by one to three decades of noise everywhere measured. Whenever the fluctuations still say something about the unitary current or the channel number, a classical fit reports a confidence interval that is too narrow. The scope of that statement has to be stated with it: the classical arm here is least squares on the deterministic mean current with the unitary current held fixed, in its four-parameter configuration. It is not the variance-mean regression, this paper does not benchmark non-stationary fluctuation analysis, and nothing here should be read as a measured comparison against it.

The map also fixes where a parameter stops being reachable at any noise level. Minimising the corrected standard error over the whole noise sweep, the channel number is unreachable below about $17$ channels and the opening rate below about $52$, at this acquisition interval. % src: projects/eLife_2025/figures/paper_both/figure_6.html (best achievable SE over the noise sweep; N_ch below which the parameter is unreachable)
The unitary current and the closing rate have no such floor within the simulated range.

The evidence is thinner than the picture, and the subsection says so rather than leaving it to a caption. The two boundaries at which least squares becomes honest, and the closure floor, are each fitted through four measured crossings. The channel count at which region~3 opens spans $120$ to $1368$ depending only on whether the strict or the loose criterion is used. % src: projects/eLife_2025/figures/paper_both/figure_6.html (region 3 OPENS at N_ch = 120-1368, the range being the two criteria)
The lower vertex, where the map closes from below, is extrapolated to $N_{\mathrm{ch}}$ between $3$ and $9$, below the simulated floor of ten channels, and the ten cells that would pin it have not been run. % src: projects/eLife_2025/figures/paper_both/figure_6.html (map CLOSES from below at N_ch = 3-9, EXTRAPOLATED)
All the boundaries are drawn at one acquisition interval, with the ribbons showing where the same boundary moves as the interval is swept from $0.01$ to $1$ in units of $\tau$. The map is a concept map, not a phase diagram. Its boundaries are level sets of continuous diagnostics, so a looser criterion moves each of them by up to a decade in noise without changing the layout.

Three recording configurations are placed on the same plane, as footprints rather than points: an excised outside-out patch, a whole cell, and a two-electrode voltage-clamped oocyte. Each footprint is the range of channel numbers the configuration expresses crossed with the range of instrumental noise it carries, converted into the sweep's own axis through the unitary current and the channel time constant, and clipped to the relaxation each configuration can resolve. In the sweep's label units the patch spans about $3\times10^{-5}$ to $2.5$, the whole cell about $10^{-4}$ to $20$, and the oocyte about $8\times10^{2}$ to $9\times10^{6}$. % src: projects/eLife_2025/figures/paper_both/figure_6.html ("three preparations, in the sweep's own label" table)
The placement is provisional and is offered as orientation. Typical instrumental noise per configuration is not published in a form that can be used directly, so the noise ranges are read off published current traces by a calibrated procedure, and the minimum resolvable current for two of the three configurations is a working value that has not been sourced. % src: docs/bibliography/recording_configurations/TRACE_WIDTH_METHOD.md (noise read off traces; calibration and its biases) ; projects/eLife_2025/figures/paper_both/figure_6.Rmd:411-413 (Imin for patch and whole cell are unsourced placeholders)

\begin{figure}[htbp]
\begin{fullwidth}
\centering
\includegraphics[width=\linewidth]{Figure_6}
\end{fullwidth}
\caption{\textbf{What a macroscopic recording can be asked for: a usage map over channel number and instrumental noise.} The plane is the design space, the number of channels in the patch against the instrumental noise label, both on logarithmic scales, at one acquisition interval $\Delta\!\cdot\!k_{\mathrm{off}}=0.1$. Four boundaries divide it, and they are two questions asked of two methods. Colour is the parameter the boundary is about. Line type is the question: solid asks whether the parameter can be measured at all, drawn where the distortion-corrected standard error crosses a factor of two; dotted asks whether the reported error bar of least squares is honest, drawn where its information distortion crosses $1.15$; dashed asks the same of the likelihood's own error bar, and is the boundary below which the Gaussian closure is misspecified. Ribbons are the same boundary swept over acquisition intervals from $0.01$ to $1$ in units of $\tau$, so a ribbon says where a reader's own sampling rate puts the boundary and is not an error bar. Dots are the measured crossings and the line through them is their fitted power law, drawn faint where it leaves the simulated range, which is marked. Vertical strips mark the two vertices, each a range set by the strict and the loose criterion. The opening rate binds before the closing rate because a single concentration jump determines it only through the difference between the two relaxation rates. The map is a concept map, not a phase diagram. Its boundaries are level sets of continuous diagnostics, so a looser criterion moves each of them by up to a decade in noise without changing the layout.} % src: projects/eLife_2025/figures/paper_both/Figure_6_caption.md and figure_6.Rmd header (encoding: one visual channel per meaning; ribbon = interval sweep)
\label{fig:usage}
% Figure 6 carries no supplements. The `regions` and `frontiers` variants were superseded on
% 2026-07-31, when the canonical map became Figure_6 plain, and figure_6.Rmd now writes them into
% figures/archive/figure_6_superseded_20260731/ rather than beside the body figure. Their
% \figsupp declarations pointed at files that no longer exist beside the manuscript and were the
% compile failure. Nothing was lost with them: the "concept map, not a phase diagram" caveat both
% carried is already the last sentence of the caption above. What is not carried anywhere now is
% the five named regions and the three recording footprints; if that is wanted back it is a new
% figure and not a restored reference. A parent at zero supplements is the eLife norm (46.6%).
\end{figure}
